

% 第一章 绪论
\chapter{绪论}
\section{研究背景与意义}
汽车工业的发展趋势正向智能化、网联化的方向发展，车辆的开发过程中仿真是一项非常重要的工作，在车辆的研发中仿真测试工作的比重越来越大。传统的车辆多体物理仿真软件有ADAMS、SIMPACK等等\cite{gao2026suspension}，它们能够完成非常高的精度求解工作，但是由于自身存在着很多难以克服的问题，如实时性不好、可视化效果差、不能实现流畅的交互等问题，这些都使得它们很难满足自动驾驶场景下的测试、底盘性能动态检验等新的研发需要

虚幻引擎是一款目前常用的实时渲染引擎，可以产生高品质的场景画面，并且能够方便地创建人机交互的功能，所以受到许多人的青睐；Chrono为开源的工程级多体动力学引擎，在精准搭建车辆各子系统的模型基础上还可以耦合求解多个物理场，与商用级别的仿真软件相比它的精度接近\cite{wang2026auv}。研究人员把Chrono所拥有的高精度物理求解的能力，同虚幻引擎所具有的高质量可视化功能结合起来使用起来，在此基础上构建出了一种具有“物理核心+可视界面”特征的车辆仿真的系统。该类系统能够很好地克服以往的仿真系统存在的一些不足，并且契合现代汽车开发朝着更加数码化以及实时化的趋势发展这一情况。

在这样一个行业中，本文主要要做的就是开发一个基于Chrono、虚幻引擎的多体物理车辆仿真的软件系统，研究的主要目的就是要解决几个关键技术问题，即多个物体的动力学模型与可视化场景之间的实时同步，汽车核心部件的准确建模方法，以及复杂情况下如何保持仿真结果的准确性等问题，这样一种模拟环境不仅适用于汽车的研发阶段工作，也可以用于自动驾驶技术的验证测试等，并具有较为显著的价值\cite{chai2026trajectory}。

\section{国内外研究现状}

\subsection{国外研究现状}

国外对于车辆多体仿真和实时可视化的融合研究比国内要早一些，在取得了一定的阶段性成果之后就停止了下去。关于多体动力学引擎的使用情况，Cherno项目组与密歇根大学等科研单位共同创建，利用chrono：vehicle模块建立并检验出诸多类型的汽车模拟器，在其中就包括对悬架、轮胎等诸多关键子系统的精确复现，之后又对其进行过严格的验证工作以证明其能够在越野车通过性的仿真中发挥预期作用，与此同时该团队还在虚拟现实软件 Unreal Engine 上成功移植了基于物理的SimCenter 3D工具集，并对其应用做了详尽分析\cite{li2026carfollowing}；另外一些企业如通用、宝马等也选择采用这种融合式的设计思路——同时依靠两个强大的可视化引擎（专业的实体级计算机模拟系统与即时渲染型视觉化编程语言）协同工作来打造数字孪生体系，在此基础之上他们还特意针对不同应用场景（底盘性能校正、自动驾驶功能评价等等）进行了具体的操作步骤说明以及详细的技术流程介绍，并且最终达成共识称这样的技术框架确实经得起时间考验并在多个方向得到了广泛应用。

除此之外，国外学者在实时数据同步领域也进行了深入的研究，他们提出了基于UDP 协议的轻量化数据传输方法，通过这种方法实现了物理仿真和可视化场景的毫秒在多物理场耦合上，他结合了Chrono和OpenFOAM使用起来非常方便，在对车辆气动阻力以及动力学性能进行了联合仿真的过程中也得到了很好的效果。但是国外的相关研究大多是为了高端车企定制开发出来的技术方案，开源的并且可以重复使用的较少，在满足通用化和低成本研发的要求上存在一定的难度。

\subsection{国内研究现状}

国内对于车辆仿真技术的研究，大多数都是把传统的商用软件（如 ADAMS）拿来应用或者二次开发，并没有针对汽车工业发展而来的全新研究，主要分为三个方面，一是基于 ADAMS 的车辆底盘性能仿真，二是基于 MATLAB / Simulink 的控制算法验证等。近数年来，随着虚幻引擎、Chrono等开源工具的逐步推广，在国内高校及科研院所中也出现了有关虚拟现实和物理模拟技术相结合的研究：清华大学利用Chrono建立电动车辆多体动力学模型并实现动力传动系统性能仿真的工作已经开展起来；吉林大学把虚幻引擎与PhysX物理引擎结合起来创建自动驾驶环境仿真实验室，并且对基本交通场景进行了可视化的测试。

但是国内的研究还有许多明显的不足之处，第一是大多数研究使用的是游戏级别的物理引擎如PhysX这类的引擎，在多体动力学求解的精度上并不高，不能满足工程级的需求；第二是在Chrono和虚幻引擎之间没有建立更深入的配合机制，对于同步延迟补偿、动力学参数校准等也没有做过系统的探索；第三就是缺少一个将场景模型同车辆物理行为有机地融合在一起的完整系统。因此，开发出一个高精度、高实时性、高耦合度的基于Chrono和虚幻引擎的车辆仿真是很有研究价值的。

\section{研究内容与方法}

\subsection{研究内容}

本研究旨在开发一套基于Chrono与虚幻引擎的多体物理车辆仿真系统，核心研究内容包括以下四个方面：

\begin{enumerate}[label=(\arabic{enumi})]
	\item 高精度车辆多体动力学模型的创建要依靠Chrono::Vehicle模块来完成车架、悬架、转向、传动、轮胎等主要子系统建立之后再加以整合的工作，在这里动力学响应准确程度要按照实际车辆测试数据5来进行参照\cite{song2026fuzzy}。
	\item 多体仿真的可视化场景之间的合作关系架构是采用物理层、数据传递层和虚拟化可视化层来实现的，在这两个可视化情景之间使用了UDP/TCP进行双向通信，并且还加入了大约百分之六的小量同步误差修正机制以减小误差积累的过程\cite{luo2026sar}。
	\item 使用虚幻引擎构建出多种类型的实验环境，铺设的试验路段、潮湿的试验路段、越野地等都是比较常见的类型，同时也有典型的速度工况以及急刹工况、蛇形转弯等复杂的路况表现形式，在此基础上来实现各种不同的仿真情况以及被测试的情况所覆盖程度\cite{gao2025pitch}。
	\item 系统功能集成和优化，创建交互界面改善Chrono求解器的计算效率、修改虚拟引擎渲染效果，在长时间的模拟之中也能够保持稳定的状态而不会出现数据丢失或者是状态混乱的现象发生的情况。
\end{enumerate}

\subsection{研究方法}
本研究采用的技术方法如下：
\begin{enumerate}[label=(\arabic{enumi})]
	\item 多体动力学建模的方法就是利用Chrono::Vehicle模块标准接口进行标准化的建立各个子系统并使它们之间产生联系的办法。
	\item 实时同步就是利用UDP协议来实现轻量级的数据交换，配合上预测迭代的方法来补偿由于延迟产生的差异，从而使得物理仿真和可视化的同步。
	\item 场景建模和渲染优化就是借助虚幻引擎的地形编辑工具、材质系统以及LOD(细节层次)功能来保持一定的环境真实性同时又不造成过高的运行消耗。
	\item 验证和改进的方法：用实际的车辆测试数据来对系统进行精度的检验，利用典型的工作环境下的模拟结果来进行比较，然后根据这些情况来修改求解器的参数以及渲染所使用的设置，并且提高系统的整体性能达到百分之八以上的效果\cite{wang2025longitudinal}。
\end{enumerate}

\section{研究创新点}
本研究的创新点主要体现在以下三个方面：
\begin{enumerate}[label=(\arabic{enumi})]
	\item 高精度物理内核与高保真的可视化外壳之间存在着协同结构，这由研究人员提出并且实施，并且该方案能够把 Chrono 能够实现的工程级多体动力学计算能力同虚幻引擎所支持下的实时光渲染相结合，而实际操作中却很少见这样的做法，但是这个协同结构可以解决上述不足之处。
	\item 一个多子系统精准耦合的车辆动力学模型是由研究人员构建起来的，对于悬架、轮胎等核心子系统，采用工程级的模型Pac2002轮毂模数，并且加入了土壤力学模块来适应复杂的地形，这样就提高了仿真的工程实用性。
	\item 研究人员提出一种低延迟数据同步的方法，在其中加入预测迭代法来补偿传输过程中产生的延时，使得物理仿真与可视化可以达到毫秒级别同步的效果，最好的情况下同步时间能达到10ms左右，保证了仿真的准确性以及一致性。
\end{enumerate}
\section{论文结构}
本文的研究工作主要集中在CarlaUE4仿真实验平台上，结合Chrono仿真算法对多体物理仿真模型进行设计、优化及实验验证，最后完成对研究成果的归纳总结，形成了完整的系统化研究体系，共有五个章节的内容构成整个论文；论文从理论基础知识入手，再经过模型的设计与完善、性能优化以及最终成果的总结三个步骤来进行相关研究，在图\ref{1.png}中可以清晰地看到本文的论文结构。

绪论是第一部分，在这部分中，研究人员对多体物理仿真技术的研究背景、应用场景和创新点作了详细的论述，并整理出国内国外有关CarlaUE4与Chrono联合仿真的研究情况，明确本文要解决的问题及全文的结构。

理论框架是第二章的主要内容，这一章分析了 CarlaUE4 仿真平台的核心特点，说明了 Chrono 多体动力学的求解原理，梳理了多体物理仿真的核心技术和优化方法，还提出了一个多平台间互相联系、动力学求解更精细的方法。

第三章主要是对模型进行改进和优化，主要研究了用 Chrono 算法实现的多体动力学求解优化、使用 CarlaUE4 和 Chrono 的耦合优化等内容，并且对动力学方程进行了改进、约束问题得到了解决、数据交互接口也设计完成，从而完成了对多体物理仿真模型的构建。

第四章为实验验证，本研究以多组对比试验、消融试验和多场景仿真的方式对模型进行改进的效果做了量化化的分析，在多个环境下多次实验之后得到的结果证明该模型具有较好的工程应用性。

本章对多体物理仿真模型的建立与优化研究结果进行了总结，并从仿真模型存在的问题出发，提出了仿真实现的方法和思路，最后阐述了在未来的发展过程中，对于该类模型进行改进的方向并给出了一些扩展研究的建议。
\begin{figure}[htbp] % 可以是h（here），t（top），b（bottom），p（page of floats）
	\centering
	\includegraphics[width=0.75\textwidth]{1} % 假设图片文件名为car.pdf或car.png等，位于当前工作目录
	\caption{论文结构} % 图片标题
	\label{1.png} % 用于引用的标签
\end{figure}

